\documentclass[12pt,a4paper]{article}

% Packages
\usepackage[a4paper, margin=2.5cm, headheight=16pt]{geometry}
\usepackage[T1]{fontenc}
\usepackage[utf8]{inputenc}
\usepackage{mathpazo}          % Palatino font for text + math
\usepackage{amsmath}
\usepackage{amssymb}
\usepackage{booktabs}
\usepackage{array}
\usepackage{tabularx}
\usepackage{xcolor}
\usepackage{colortbl}
\usepackage{titlesec}
\usepackage{enumitem}
\usepackage{fancyhdr}
\usepackage{mdframed}
\usepackage{parskip}
\usepackage{hyperref}

% Global table row height & cell padding
\setlength{\extrarowheight}{6pt}
\setlength{\tabcolsep}{12pt}

% Colors
\definecolor{headerblue}{RGB}{30, 80, 160}
\definecolor{lightblue}{RGB}{220, 235, 250}
\definecolor{questionbg}{RGB}{240, 248, 255}
\definecolor{answerbg}{RGB}{250, 252, 255}
\definecolor{tablehead}{RGB}{50, 100, 180}

% Page style
\pagestyle{fancy}
\fancyhf{}
\fancyhead[L]{\textcolor{headerblue}{\textbf{Machine Learning Theory}}}
\fancyhead[R]{\textcolor{headerblue}{\textit{Exam Questions \& Answers}}}
\fancyfoot[C]{\textcolor{gray}{\thepage}}
\renewcommand{\headrulewidth}{1.5pt}
\renewcommand{\headrule}{\hbox to\headwidth{\color{headerblue}\leaders\hrule height \headrulewidth\hfill}}

% Question box style
\mdfdefinestyle{questionstyle}{
    backgroundcolor=questionbg,
    linecolor=headerblue,
    linewidth=1.5pt,
    leftline=true,
    rightline=false,
    topline=false,
    bottomline=false,
    innerleftmargin=10pt,
    innerrightmargin=10pt,
    innertopmargin=6pt,
    innerbottommargin=6pt,
}

% Answer box style
\mdfdefinestyle{answerstyle}{
    backgroundcolor=white,
    linecolor=gray!30,
    linewidth=0.5pt,
    innerleftmargin=10pt,
    innerrightmargin=10pt,
    innertopmargin=6pt,
    innerbottommargin=6pt,
}

% Table spacing
\renewcommand{\arraystretch}{1.7}

% Table column types
\newcolumntype{H}{>{\columncolor{tablehead}\color{white}\bfseries}l}
\newcolumntype{L}{>{\raggedright\arraybackslash}X}

% Custom commands
\newcommand{\question}[2]{%
    \begin{mdframed}[style=questionstyle]
        \textbf{\textcolor{headerblue}{Q#1.}} \textbf{#2}
    \end{mdframed}%
    \vspace{2pt}%
}

\newcommand{\ans}{\textbf{\textcolor{headerblue}{Answer:}}\\[4pt]}
\newcommand{\subhead}[1]{\textbf{\textcolor{headerblue}{#1}}}

\begin{document}

% Title Page
\begin{titlepage}
    \centering
    \vspace*{2cm}
    {\Huge\bfseries\textcolor{headerblue}{Machine Learning}\par}
    \vspace{0.5cm}
    {\LARGE\textcolor{gray}{Exam Questions with Answers}\par}
    \vspace{1.5cm}
    \rule{\textwidth}{1.5pt}
    \vspace{0.3cm}
    \rule{\textwidth}{0.5pt}
    \vspace{2cm}
    {\large Prepared by: \textbf{Sharika}\par}
    \vfill
    {\small All answers are written in simple English for easy understanding.}
\end{titlepage}

\setlength{\parskip}{6pt}

%===============================================================
\question{1}{Describe three methods for handling missing data and explain when each would be appropriate.}

\ans

\subhead{Method 1: Deletion (Removing Missing Data)}

We remove rows or columns that have missing values.

\textit{When to use:} Use this when very few rows have missing data (less than 5\%). If you remove too many rows, you lose important information.

\subhead{Method 2: Mean / Median / Mode Imputation (Filling with a Value)}

We replace missing values with the average (mean), middle value (median), or most common value (mode) of that column.

\textit{When to use:} Use mean for numbers like age or salary. Use median when data has extreme values. Use mode for categories like gender or color.

\subhead{Method 3: Prediction-Based Imputation (Using a Model to Fill)}

We use a machine learning model to predict what the missing value should be, based on other data.

\textit{When to use:} Use this when many values are missing and accuracy is very important. It is more complex but gives better results.

\vspace{10pt}
%===============================================================
\question{2}{Define cross-validation. Explain the purpose of using K-fold cross-validation and describe the steps involved in applying it to a dataset.}

\ans

Cross-validation is a method to test how well a machine learning model works on new data. It helps us check if the model is learning well or just memorizing the training data.

\subhead{What is K-Fold Cross-Validation?}

K-fold cross-validation splits the data into $K$ equal parts (called folds). The model is trained and tested $K$ times. Each time, a different fold is used for testing and the rest for training.

\subhead{Purpose:}
\begin{itemize}[leftmargin=*]
    \item It gives a more reliable score for the model.
    \item It uses all data for both training and testing.
    \item It helps detect overfitting (when the model is too focused on training data).
\end{itemize}

\subhead{Steps:}
\begin{enumerate}[leftmargin=*]
    \item Split the dataset into $K$ equal parts (e.g., $K = 5$).
    \item Use fold 1 as the test set. Use folds 2, 3, 4, 5 for training.
    \item Train the model and record the accuracy.
    \item Repeat with fold 2 as the test set, and so on.
    \item Calculate the average accuracy from all $K$ rounds.
\end{enumerate}

\vspace{10pt}
%===============================================================
\question{3}{Compare between Bagging and Boosting algorithm.}

\ans

\begin{tabularx}{\textwidth}{|>{\bfseries}l|L|L|}
\hline
\rowcolor{tablehead}\textcolor{white}{\textbf{Feature}} & \textcolor{white}{\textbf{Bagging}} & \textcolor{white}{\textbf{Boosting}} \\
\hline
Training & Models train at the same time (parallel) & Models train one after another (sequential) \\
\hline
Focus & Each model gets equal importance & Each new model focuses on the mistakes of the last one \\
\hline
Goal & Reduces variance (overfitting) & Reduces bias (underfitting) \\
\hline
Example & Random Forest & AdaBoost, XGBoost \\
\hline
Speed & Faster (parallel) & Slower (sequential) \\
\hline
\end{tabularx}

\vspace{10pt}
%===============================================================
\question{4}{What is Supervised Machine Learning? Explain how does it work.}

\ans

Supervised machine learning is a type of machine learning where the computer learns from \textbf{labeled data}. Labeled data means each input has a correct answer (label) already given.

\subhead{How Does It Work?}
\begin{enumerate}[leftmargin=*]
    \item \textbf{Give labeled data:} We give the model many examples. For example: house size = input, house price = label.
    \item \textbf{Model learns:} The model finds patterns between the input and the label.
    \item \textbf{Make predictions:} After learning, the model can predict the answer for new inputs it has not seen before.
    \item \textbf{Check accuracy:} We compare the model's predictions with the real answers to see how well it learned.
\end{enumerate}

\textit{Examples:} Predicting house prices (regression), classifying emails as spam or not spam (classification).

\vspace{10pt}
%===============================================================
\question{5}{Explain the concept of k-fold cross-validation and its advantages over simple train-test splits in case of imbalanced dataset.}

\ans

K-fold cross-validation divides the data into $K$ equal groups. The model is trained $K$ times, using a different group for testing each time. The average score from all runs gives a final result.

\subhead{Advantages over Simple Train-Test Split (Especially for Imbalanced Datasets):}
\begin{itemize}[leftmargin=*]
    \item \textbf{Better use of data:} Every data point is used for both training and testing.
    \item \textbf{More reliable result:} The average of $K$ scores is more stable than one single score.
    \item \textbf{Handles imbalance better:} In an imbalanced dataset (e.g., 90\% class A, 10\% class B), a simple split may put all minority class examples in training only. K-fold ensures the minority class appears in each test fold.
    \item \textbf{Less bias:} It reduces the chance of getting lucky (or unlucky) with one particular split.
\end{itemize}

\vspace{10pt}
%===============================================================
\question{6}{List down five differences between Machine Learning and Artificial Intelligence.}

\ans

\begin{tabularx}{\textwidth}{|c|L|L|}
\hline
\rowcolor{tablehead}\textcolor{white}{\textbf{\#}} & \textcolor{white}{\textbf{Artificial Intelligence (AI)}} & \textcolor{white}{\textbf{Machine Learning (ML)}} \\
\hline
1 & AI is the big field of making smart machines. & ML is a part of AI that focuses on learning from data. \\
\hline
2 & AI can use rules, logic, and reasoning. & ML only learns from examples and patterns. \\
\hline
3 & AI does not always need data to work. & ML always needs data to learn. \\
\hline
4 & AI tries to copy human thinking. & ML tries to improve predictions automatically. \\
\hline
5 & Examples: robots, chess programs. & Examples: spam filters, recommendation systems. \\
\hline
\end{tabularx}

\vspace{10pt}
%===============================================================
\question{7}{Write down the differences between Supervised and Unsupervised Machine Learning.}

\ans

\begin{tabularx}{\textwidth}{|>{\bfseries}l|L|L|}
\hline
\rowcolor{tablehead}\textcolor{white}{\textbf{Feature}} & \textcolor{white}{\textbf{Supervised Learning}} & \textcolor{white}{\textbf{Unsupervised Learning}} \\
\hline
Labels & Data has labels (answers given) & Data has no labels \\
\hline
Goal & Predict an output & Find hidden patterns or groups \\
\hline
Example & Spam detection, price prediction & Customer grouping, topic detection \\
\hline
Feedback & Model learns from correct answers & No correct answers given \\
\hline
Algorithms & Linear Regression, SVM & K-Means, DBSCAN \\
\hline
\end{tabularx}

\vspace{10pt}
%===============================================================
\question{8}{Describe Hard Margin of SVM.}

\ans

SVM stands for \textbf{Support Vector Machine}. It is a classification algorithm that draws a line (or boundary) to separate two classes of data.

\subhead{What is Hard Margin SVM?}

Hard margin SVM finds the best line (called the \textit{decision boundary} or \textit{hyperplane}) that separates two classes with the largest possible gap (margin). It assumes the data is perfectly separable --- meaning no data point crosses to the wrong side.

\subhead{Key Terms:}
\begin{itemize}[leftmargin=*]
    \item \textbf{Decision Boundary:} The line that separates the two classes.
    \item \textbf{Margin:} The distance between the decision boundary and the nearest data points from each class.
    \item \textbf{Support Vectors:} The data points that are closest to the decision boundary. These points define the margin.
\end{itemize}

\subhead{Goal:}

Maximize the margin between the two classes. The formula is:
\[
\text{Maximize} \quad \frac{2}{\|\mathbf{w}\|}
\]
where $\mathbf{w}$ is the weight vector that defines the boundary.

\subhead{Limitation:}

Hard margin SVM only works when the data is perfectly separable (no overlapping points). Real-world data often has noise, so Soft Margin SVM is used instead.

\vspace{10pt}
%===============================================================
\question{9}{Describe briefly about univariate linear regression in machine learning.}

\ans

Univariate linear regression is the simplest type of regression. It uses \textbf{one input} (feature) to predict \textbf{one output} (label).

\textit{Example:} We use house size ($X$) to predict house price ($Y$).

\subhead{Hypothesis Function:}
\[
h(x) = \theta_0 + \theta_1 x
\]
Here, $\theta_0$ is the intercept (starting point) and $\theta_1$ is the slope (how much $Y$ changes when $X$ increases by 1).

\subhead{Cost Function (How to Measure Accuracy):}
\[
J(\theta_0, \theta_1) = \frac{1}{2m} \sum_{i=1}^{m} \left(h(x^{(i)}) - y^{(i)}\right)^2
\]
This calculates the average squared difference between predicted and actual values. A smaller value means the model is more accurate.

\vspace{10pt}
%===============================================================
\question{10}{Briefly describe the gradient descent algorithm for linear regression.}

\ans

Gradient descent is a method to find the best weights $(w_0, w_1, \ldots, w_n)$ that minimize the cost function. It repeats small steps to reduce the error.

\subhead{Algorithm: \texttt{GRADIENTDESCENT}$(X, y, \alpha, \text{maxIter})$}

\begin{enumerate}[leftmargin=*]
    \item \textbf{for} $j = 1$ to $m$: set $x_0^{(j)} = 1$ \hfill \textit{[Add intercept feature to each example]}
    \item Initialize $w_0, w_1, \ldots, w_n$ randomly
    \item $\text{iter} = 0$
    \item \textbf{while} $\text{iter++} \leq \text{maxIter}$:
    \begin{enumerate}
        \item \textbf{for} $j = 0$ to $n$: $\text{slope}_j = 0$ \hfill \textit{[Reset slopes to zero]}
        \item \textbf{for} $i = 1$ to $m$: \hfill \textit{[Loop over all training examples]}
        \begin{enumerate}
            \item $\hat{y} = w_0 + w_1 x_1^{(i)} + \cdots + w_n x_n^{(i)}$ \hfill \textit{[Predict]}
            \item $e = \hat{y} - y^{(i)}$ \hfill \textit{[Calculate error]}
            \item \textbf{for} $j = 0$ to $n$: $\text{slope}_j = \text{slope}_j + e \cdot x_j^{(i)}$ \hfill \textit{[Accumulate slopes]}
        \end{enumerate}
        \item \textbf{for} $j = 0$ to $n$: $w_j = w_j - \alpha \cdot \text{slope}_j$ \hfill \textit{[Update weights]}
    \end{enumerate}
    \item \textbf{return} $w_0, w_1, \ldots, w_n$
\end{enumerate}

\subhead{Key Terms:}
\begin{itemize}[leftmargin=*]
    \item $\alpha$ = learning rate. Controls the step size.
    \item $\text{slope}_j$ = total gradient for weight $j$.
    \item $e$ = error = predicted value minus actual value.
\end{itemize}

\vspace{10pt}
%===============================================================
\question{11}{How gradient descent automatically improves the hypothesis function.}

\ans

Gradient descent automatically updates $\theta_0$ and $\theta_1$ in each step. After each update, the hypothesis function $h(x) = \theta_0 + \theta_1 x$ gets closer to the real data.

As $\theta$ values change, the line moves and fits the data better. The cost $J$ decreases with each step. When the cost stops decreasing, the model has found the best fit line.

This is how gradient descent automatically improves the model without us manually choosing the slope and intercept.

\vspace{10pt}
%===============================================================
\question{12}{What is a classification problem in machine learning? Give an example.}

\ans

A classification problem is when the output (label) belongs to a \textbf{specific category}, not a number. The model learns to place inputs into one of these categories.

\subhead{Example:}

\textit{Email classification:} The input is the email text. The output is either ``spam'' or ``not spam''. This is a \textbf{binary classification} (two categories).

\subhead{Other Examples:}
\begin{itemize}[leftmargin=*]
    \item Predicting if a tumor is malignant (cancer) or benign (not cancer).
    \item Recognizing handwritten digits (0 to 9) --- multi-class classification.
    \item Classifying a photo as cat, dog, or bird.
\end{itemize}

\vspace{10pt}
%===============================================================
\question{13}{Describe briefly the logistic regression model and cost function of a single example.}

\ans

\subhead{Logistic Regression Model:}

Logistic regression is used for classification. Unlike linear regression which gives any number, logistic regression gives a probability between 0 and 1.

\textbf{Hypothesis (Sigmoid function):}
\[
h(x) = \frac{1}{1 + e^{-\boldsymbol{\theta}^T \mathbf{x}}}
\]

If $h(x) \geq 0.5$, we predict class 1. If $h(x) < 0.5$, we predict class 0.

\subhead{Cost Function for a Single Example:}
\[
\text{Cost}(h(x), y) = \begin{cases}
-\log(h(x)) & \text{if } y = 1 \\
-\log(1 - h(x)) & \text{if } y = 0
\end{cases}
\]

If $y = 1$ and our prediction $h(x)$ is close to 1, cost is low (good). If our prediction is far from 1, cost is high (bad). This pushes the model to predict correctly.

\vspace{10pt}
%===============================================================
\question{14}{Briefly describe the process to develop a gradient descent for linear regression with regularization.}

\ans

Regularization is a technique to prevent overfitting. It adds a penalty to the cost function for large $\theta$ values.

\subhead{Regularized Cost Function (Ridge / L2):}
\[
J(\boldsymbol{\theta}) = \frac{1}{2m} \sum_{i=1}^{m}\left(h(x^{(i)}) - y^{(i)}\right)^2 + \frac{\lambda}{2m} \sum_{j=1}^{n} \theta_j^2
\]

Here, $\lambda$ is the regularization parameter. It controls how much we penalize large $\theta$ values.

\subhead{Gradient Descent Update with Regularization:}
\[
\theta_0 := \theta_0 - \alpha \cdot \frac{1}{m} \sum_{i=1}^{m}\left(h(x^{(i)}) - y^{(i)}\right) x_0^{(i)} \quad \text{[No regularization for } \theta_0\text{]}
\]
\[
\theta_j := \theta_j \left(1 - \alpha\frac{\lambda}{m}\right) - \alpha \cdot \frac{1}{m} \sum_{i=1}^{m}\left(h(x^{(i)}) - y^{(i)}\right) x_j^{(i)} \quad \text{for } j \geq 1
\]

The term $\left(1 - \alpha\frac{\lambda}{m}\right)$ slightly shrinks $\theta_j$ in each step, which prevents it from becoming too large.

\vspace{10pt}
%===============================================================
\question{15}{Derive the algorithm of Support Vector Machine (SVM). What is the use of an SVM?}

\ans

\subhead{What is SVM?}

SVM is a supervised learning algorithm used for classification. It finds the best boundary (hyperplane) that separates two classes with the maximum margin.

\subhead{SVM Algorithm (Derivation):}
\begin{itemize}[leftmargin=*]
    \item Given data points $(x^{(i)}, y^{(i)})$ where $y^{(i)} \in \{+1, -1\}$.
    \item We want to find $\mathbf{w}$ (weights) and $b$ (bias) such that:
    \[
    \mathbf{w}^T \mathbf{x} + b \geq 1 \quad \text{for class } +1
    \]
    \[
    \mathbf{w}^T \mathbf{x} + b \leq -1 \quad \text{for class } -1
    \]
    \item The margin between the two classes $= \dfrac{2}{\|\mathbf{w}\|}$.
    \item To maximize the margin, we minimize $\dfrac{\|\mathbf{w}\|^2}{2}$.
    \item \textbf{Optimization:} Minimize $\dfrac{1}{2}\|\mathbf{w}\|^2$ subject to $y^{(i)}\left(\mathbf{w}^T \mathbf{x}^{(i)} + b\right) \geq 1$.
\end{itemize}

\subhead{Uses of SVM:}
\begin{itemize}[leftmargin=*]
    \item Email spam detection.
    \item Image classification (e.g., face recognition).
    \item Medical diagnosis (cancer detection).
    \item Text classification (news categories).
\end{itemize}

\vspace{10pt}
%===============================================================
\question{16}{Suppose you have $m=8$ training examples with $n=3$ features. The normal equation is $\boldsymbol{\theta} = (X^T X)^{-1} X^T \mathbf{y}$. What are the dimensions of $\boldsymbol{\theta}$, $X$, and $\mathbf{y}$?}

\ans

$m = 8$ training examples, $n = 3$ features. We add 1 extra feature ($x_0 = 1$) for the intercept, so total features $= n + 1 = 4$.

\subhead{Dimensions:}
\begin{itemize}[leftmargin=*]
    \item $X$: has $m$ rows and $(n+1)$ columns. So $X$ is $8 \times 4$.
    \item $\mathbf{y}$: has $m$ rows and 1 column. So $\mathbf{y}$ is $8 \times 1$.
    \item $\boldsymbol{\theta}$: has $(n+1)$ rows and 1 column. So $\boldsymbol{\theta}$ is $4 \times 1$.
\end{itemize}

\subhead{Verification:}

$X^T$ is $4 \times 8$. \quad $X^T X$ is $4 \times 4$. \quad $(X^T X)^{-1}$ is $4 \times 4$. \quad $X^T \mathbf{y}$ is $4 \times 1$.

So $\boldsymbol{\theta} = (X^T X)^{-1} X^T \mathbf{y}$ is $4 \times 1$. \checkmark

\vspace{10pt}
%===============================================================
\question{17}{In the context of supervised learning, what distinguishes the training data from the testing data, and why is this differentiation important?}

\ans

\subhead{Training Data:}

This is the data we use to \textbf{teach} the model. The model sees the inputs and labels and learns patterns from them.

\subhead{Testing Data:}

This is data the model has \textbf{never seen} before. We use it to check how well the model predicts on new examples.

\subhead{Why is this Differentiation Important?}
\begin{itemize}[leftmargin=*]
    \item It checks if the model actually learned or just memorized the training data (overfitting).
    \item It gives a real-world measure of model performance.
    \item It prevents us from reporting false accuracy (a model that only works on training data is not useful).
\end{itemize}

\vspace{10pt}
%===============================================================
\question{18}{How does bagging work as a method of increasing accuracy? How does boosting compare with bagging?}

\ans

\subhead{How Bagging Works:}
\begin{itemize}[leftmargin=*]
    \item Bagging stands for \textbf{Bootstrap Aggregating}.
    \item It creates many small datasets by randomly sampling from the original data (with replacement).
    \item It trains a separate model on each small dataset at the same time (parallel).
    \item Final prediction is made by \textbf{voting} (for classification) or \textbf{averaging} (for regression).
    \item Bagging reduces \textbf{variance} --- it makes the model more stable and less likely to overfit.
\end{itemize}

\subhead{How Boosting Compares:}
\begin{itemize}[leftmargin=*]
    \item Boosting trains models one after another (sequential).
    \item Each new model focuses more on the examples that the previous model got wrong.
    \item Boosting reduces \textbf{bias} --- it makes weak models stronger over time.
    \item Boosting can overfit if not controlled (e.g., with regularization or early stopping).
\end{itemize}

\vspace{10pt}
%===============================================================
\question{19}{Explain the concept of a labeled dataset and its significance in the context of supervised learning algorithms.}

\ans

\subhead{What is a Labeled Dataset?}

A labeled dataset is a collection of data where each input example has a correct output (answer) already attached to it.

\subhead{Example:}
\begin{itemize}[leftmargin=*]
    \item Input: An email message. \quad Label: `Spam' or `Not Spam'.
    \item Input: House size (1500 sq ft). \quad Label: Price (\$250,000).
\end{itemize}

\subhead{Why is it Significant in Supervised Learning?}
\begin{itemize}[leftmargin=*]
    \item The model learns by comparing its predictions to the correct labels.
    \item Without labels, the model does not know if it is right or wrong.
    \item More labeled data usually means better model performance.
    \item Labeled data is the foundation of all supervised learning tasks.
\end{itemize}

\vspace{10pt}
%===============================================================
\question{20}{Write the regularized logistic regression cost function, including both the data term and the regularization term.}

\ans

The regularized logistic regression cost function is:
\[
J(\boldsymbol{\theta}) = -\frac{1}{m} \sum_{i=1}^{m} \left[ y^{(i)} \log\left(h(x^{(i)})\right) + \left(1 - y^{(i)}\right) \log\left(1 - h(x^{(i)})\right) \right] + \frac{\lambda}{2m} \sum_{j=1}^{n} \theta_j^2
\]

\begin{itemize}[leftmargin=*]
    \item \textbf{First part:} Data term --- measures how wrong the predictions are.
    \item \textbf{Second part:} Regularization term --- penalizes large $\theta$ values to prevent overfitting.
    \item \textbf{Note:} $\theta_0$ is \textbf{NOT} included in the regularization term.
\end{itemize}

\vspace{10pt}
%===============================================================
\question{21}{Derive the gradient of the cost function (regularized logistic regression).}

\ans

\subhead{For $\theta_0$ (no regularization):}
\[
\frac{\partial}{\partial \theta_0} J(\boldsymbol{\theta}) = \frac{1}{m} \sum_{i=1}^{m} \left(h(x^{(i)}) - y^{(i)}\right) x_0^{(i)}
\]

\subhead{For $\theta_j$ where $j \geq 1$ (with regularization):}
\[
\frac{\partial}{\partial \theta_j} J(\boldsymbol{\theta}) = \frac{1}{m} \sum_{i=1}^{m} \left(h(x^{(i)}) - y^{(i)}\right) x_j^{(i)} + \frac{\lambda}{m} \theta_j
\]

\subhead{Gradient Descent Update:}
\[
\theta_0 := \theta_0 - \alpha \cdot \frac{1}{m} \sum_{i=1}^{m} \left(h(x^{(i)}) - y^{(i)}\right) x_0^{(i)}
\]
\[
\theta_j := \theta_j - \alpha \left[ \frac{1}{m} \sum_{i=1}^{m} \left(h(x^{(i)}) - y^{(i)}\right) x_j^{(i)} + \frac{\lambda}{m} \theta_j \right] \quad \text{for } j \geq 1
\]

\vspace{10pt}
%===============================================================
\question{22}{Explain how bias and variance are associated with logistic regression.}

\ans

\subhead{Bias:}

Bias means the model is \textbf{too simple}. It makes wrong assumptions. High bias causes \textbf{underfitting} --- the model does not learn enough from the training data. In logistic regression, using very few features or a simple model causes high bias.

\subhead{Variance:}

Variance means the model is \textbf{too complex}. It learns too much from the training data, including noise. High variance causes \textbf{overfitting} --- the model works well on training data but badly on new data. In logistic regression, using too many features or polynomial terms causes high variance.

\subhead{Goal:}

We want \textbf{low bias} and \textbf{low variance}. This is called the \textit{bias-variance tradeoff}.

\vspace{10pt}
%===============================================================
\question{23}{Discuss how regularization techniques can impact bias and variance in logistic regression models. Provide examples to illustrate.}

\ans

Regularization adds a penalty for large $\theta$ values. It controls model complexity.

\subhead{Effect of Regularization:}
\begin{itemize}[leftmargin=*]
    \item \textbf{High $\lambda$ (strong regularization):} $\theta$ values are forced to be very small. The model becomes simpler. This \textit{reduces variance} but \textit{increases bias}. Risk: underfitting.
    \item \textbf{Low $\lambda$ (weak regularization):} Little penalty. The model can become complex. This \textit{reduces bias} but \textit{increases variance}. Risk: overfitting.
\end{itemize}

\subhead{Example:}

Suppose we classify if a tumor is cancerous.
\begin{itemize}[leftmargin=*]
    \item With $\lambda = 0$ (no regularization): the model uses complex curves to perfectly fit training data but fails on new data (\textbf{overfitting}).
    \item With $\lambda = 10$ (high regularization): the model is too simple and misses real patterns (\textbf{underfitting}).
    \item With $\lambda = 1$ (balanced): the model generalizes well to new data. \checkmark
\end{itemize}

\vspace{10pt}
%===============================================================
\question{24}{Differentiate between gradient descent and normal equation.}

\ans

\begin{tabularx}{\textwidth}{|>{\bfseries}l|L|L|}
\hline
\rowcolor{tablehead}\textcolor{white}{\textbf{Feature}} & \textcolor{white}{\textbf{Gradient Descent}} & \textcolor{white}{\textbf{Normal Equation}} \\
\hline
Approach & Iterative --- repeats many steps & Direct --- one step calculation \\
\hline
Learning Rate & Needs $\alpha$ (learning rate) & No learning rate needed \\
\hline
Speed (large data) & Works well with large datasets & Very slow for large datasets \\
\hline
Formula & $w_j := w_j - \alpha \cdot \text{slope}_j$ & $\boldsymbol{\theta} = (X^T X)^{-1} X^T \mathbf{y}$ \\
\hline
Scalability & Scales to millions of features & Fails when $n$ is very large ($> 10{,}000$) \\
\hline
\end{tabularx}

\vspace{10pt}
%===============================================================
\question{25}{Write the Revised Gradient Descent Algorithm (Vectorized form).}

\ans

The Revised Gradient Descent is a faster, vectorized version. It uses matrix operations instead of loops, making it much more efficient.

\subhead{Algorithm: \texttt{GRADIENTDESCENT}$(X, y, \alpha, \text{maxIter})$}

\begin{enumerate}[leftmargin=*]
    \item $X_E = [\mathbf{1}; X]$ \hfill \textit{[Add a column of 1s to $X$]}
    \item $\mathbf{W} = [w_0, w_1, \ldots, w_n]^T$ \hfill \textit{[Initialize weights randomly]}
    \item $\text{iter} = 0$
    \item \textbf{while} $\text{iter++} \leq \text{maxIter}$:
    \begin{enumerate}
        \item $\hat{\mathbf{Y}} = X_E \cdot \mathbf{W}$ \hfill \textit{[Predict all outputs at once]}
        \item $\mathbf{E} = \hat{\mathbf{Y}} - \mathbf{Y}$ \hfill \textit{[Calculate all errors at once]}
        \item $\mathbf{S} = X_E^T \cdot \mathbf{E}$ \hfill \textit{[Calculate slopes using transpose]}
        \item $\mathbf{W} = \mathbf{W} - \alpha \cdot \mathbf{S}$ \hfill \textit{[Update all weights at once]}
    \end{enumerate}
    \item \textbf{return} $w_0, w_1, \ldots, w_n$
\end{enumerate}

\subhead{Why is this better?}
\begin{itemize}[leftmargin=*]
    \item No inner for-loops. All calculations are done with matrices.
    \item It is much faster, especially for large datasets.
    \item $X_E$ is the extended $X$ matrix with 1s added for the bias term ($w_0$).
\end{itemize}

\vspace{10pt}
%===============================================================
\question{26}{Write the Lighter Gradient Descent Algorithm (Stochastic / Single-example update).}

\ans

The Lighter Gradient Descent updates the slope using only \textbf{ONE} training example per iteration step, instead of all $m$ examples. This makes each step faster but noisier.

\subhead{Algorithm: \texttt{LIGHTERGRADIENTDESCENT}$(X, y, \alpha, \text{maxIter})$}

\begin{enumerate}[leftmargin=*]
    \item \textbf{for} $j = 1$ to $m$: set $x_0^{(j)} = 1$ \hfill \textit{[Add intercept feature]}
    \item Initialize $w_0, w_1, \ldots, w_n$ randomly
    \item $\text{iter} = 0$
    \item \textbf{while} $\text{iter++} \leq \text{maxIter}$:
    \begin{enumerate}
        \item \textbf{for} $j = 0$ to $n$: $\text{slope}_j = 0$ \hfill \textit{[Reset slopes]}
        \item $i = \text{iter}$ \hfill \textit{[Use one example: the $i$-th example]}
        \item $\hat{y} = w_0 + w_1 x_1^{(i)} + \cdots + w_n x_n^{(i)}$ \hfill \textit{[Predict]}
        \item $e = \hat{y} - y^{(i)}$ \hfill \textit{[Calculate error for this one example]}
        \item \textbf{for} $j = 0$ to $n$: $\text{slope}_j = \text{slope}_j + e \cdot x_j^{(i)}$ \hfill \textit{[Accumulate slopes]}
        \item \textbf{for} $j = 0$ to $n$: $w_j = w_j - \alpha \cdot \text{slope}_j$ \hfill \textit{[Update weights]}
    \end{enumerate}
    \item \textbf{return} $w_0, w_1, \ldots, w_n$
\end{enumerate}

\subhead{Key Difference from Standard Gradient Descent:}
\begin{itemize}[leftmargin=*]
    \item \textbf{Standard GD:} uses ALL $m$ examples to compute slope before updating.
    \item \textbf{Lighter GD:} uses only ONE example ($i = \text{iter}$) per update step.
    \item Lighter GD is faster per step but may not converge as smoothly.
\end{itemize}

\vspace{1cm}
\begin{center}
\rule{0.5\textwidth}{0.5pt}\\[6pt]
\textit{All answers are written in simple English for easy understanding.}
\end{center}

\end{document}
